\documentclass[10pt,a4paper]{amsart}
\usepackage[margin=19mm]{geometry}
\usepackage{amsmath,amssymb,mathtools}
\usepackage[scr=rsfs,cal=euler]{mathalfa}
\usepackage{xfrac}
\usepackage[unicode,hidelinks,bookmarks=false]{hyperref}
\newcommand{\ie}{i.e.\ }
\newcommand{\cf}{cf.\ }
\newcommand{\resp}{respectively\ }
\newcommand{\Subsection}{\subsection}
\providecommand{\nobreakdash}{\nobreak\hbox{-}}
\setlength{\emergencystretch}{2em}
\multlinegap0pt
\usepackage{bm}
\usepackage{bbm}
\usepackage{dsfont}
\usepackage{xspace}
\usepackage{cancel}
\usepackage{mathtools}
\usepackage{amsthm}
\makeatletter
\@ifundefined{lemma}{\newtheorem{lemma}{Lemma}}{}
\@ifundefined{proposition}{\newtheorem{proposition}{Proposition}}{}
\makeatother
\title[Hitting probabilities of constrained random walks representi]{Hitting probabilities of constrained random walks representing tandem networks}
\author{Ali Devin Sezer}
\date{Source: arXiv:2105.05474v1 (CC BY 4.0)}
\begin{document}
\maketitle

\noindent\emph{Source and licence.} Hitting probabilities of constrained random walks representing tandem networks. Version arXiv:2105.05474v1 (CC BY 4.0), distributed under the Creative Commons Attribution 4.0 International licence, as recorded in the arXiv licence metadata for this exact version and shown on the abstract page https://arxiv.org/abs/2105.05474v1. Full rights notice: licenses/arxiv-2105.05474-CC-BY-4.0.txt.\par\smallskip

\noindent\textit{Editorial scope.} Counted unit: the Proposition whose source label is \texttt{None} in the article (source0.tex lines 1063--1065, source label \texttt{None}), delivered together with the article's own complete original proof of it (source0.tex lines 1066--1262, one \texttt{proof} environment of 197 lines). No mathematics is written, repaired or re-derived: statement and proof are the article's own text line for line. Required context retained uncounted: p:hkr (lines 794--808); d:gamma (lines 853--856); p:Ysuperharmonicfuns (lines 862--865); l:inequalityoverjumps (lines 1013--1020). Every definition, display and label of those lines is retained unchanged. Omitted from this package and not redistributed: 1--793, 809--852, 857--861, 866--1012, 1021--1062, 1263--2896. Nothing inside a retained excerpt is omitted or altered: every retained body is delivered exactly as the article's lines are written, so no omission receipt of any kind accompanies this package. Citations: the retained excerpts carry no citation command at all, so no bibliography entry of the article is transcribed and no reference is redistributed. Reference adaptation: none was needed and none was made; every reference inside the delivered unit resolves inside it, so no unresolved reference or citation is produced. Printed slips kept literal: no printed word, symbol, hypothesis or number of the counted statement or of its proof was altered. Layout and class: the article's own class is \texttt{article} with packages including geometry, amscd, mathrsfs, epsfig, amsmath, amssymb, bm, amsthm, verbatim, url, color, hyperref, makeidx, mathabx; the wrapper is the lane's pinned \texttt{amsart} editorial wrapper at 10pt on A4, which restates the theorem-like environments the retained text needs and adds the article's own macro definitions that the retained text needs (none), with extra wrapper packages (bm, bbm, dsfont, xspace, cancel, mathtools). The delivered numbering is the unit's own. Assets: no retained span contains a figure environment, a graphics-inclusion command, a PDF-page inclusion, a table or a TikZ picture; no image file is copied into this package, the package declares no asset dependency and no image file is redistributed with it. Licence: version arXiv:2105.05474v1 of this article is distributed under the Creative Commons Attribution 4.0 International licence, as recorded in the arXiv licence metadata for this exact version and shown on the abstract page https://arxiv.org/abs/2105.05474v1. A full rights notice is delivered at licenses/arxiv-2105.05474-CC-BY-4.0.txt. Characters: every retained excerpt is pure ASCII, so the delivered engine, XeLaTeX with its default Latin Modern text font, reports no missing character.
\begin{proposition}\label{p:hkr}
The function $h_{k,r}$ satisfies
\begin{equation}\label{e:howsuperharmonic}
{\mathbb E}_y[h_{k,r}(Y_1)] -h_{k,r}(y) =
\begin{cases}
h_{k,r}(y) \left(
\lambda \left(\frac{1}{r}-1\right) + \mu_1 (r-1)\right), &\text{if } k=1,\\
h_{k,r}(y) \left(
\lambda \left(\frac{1}{r}-1\right) + \mu_k (r-1) {\bm 1}_{\partial_k^c}(y)\right),
&\text{if } k \in \{2,3,...,d\}.
\end{cases}
\end{equation}
In particular, for $r \in (\rho,1)$,  $h_{1,r}$ is $Y$-superharmonic and for 
$k \in \{2,3,4,...,d\}$ $h_{k,r}$ is $Y$-superharmonic on $D_Y - \partial_k.$
\end{proposition}

\begin{equation}\label{d:gamma}
\gamma_1 \doteq 1, 
\gamma_k \doteq \frac{1}{d}\frac{\min_{j < k}\gamma_j(\lambda(1 -1/r) + \mu_{j}(1-r))}{\lambda(1/r - 1) }, k=2,3,...,d.
\end{equation}

\begin{proposition}\label{p:Ysuperharmonicfuns}
For any $k=1,2,3,...d$ and $r \in (\rho,1)$
the function $h_{2,k,r}$ is $Y$-superharmonic
\end{proposition}

\begin{lemma}\label{l:inequalityoverjumps}
For $k \in \{2,3,...,d\}$  and $r \in (\rho,1)$
\begin{equation}\label{e:lowerboundonratio}
\min_{y \in \partial_{k}} \frac{h_{2,k-1,r}(y)}{h_{2,k,r}(y)} \ge 
\gamma_{k-1,j}
\end{equation}
for $y \in \partial_{k}.$
\end{lemma}

\begin{proposition}
The process $\{S_k, k=0,1,2,3,...\}$ is a supermartingale.
\end{proposition}

\begin{proof}
The proof is a case by case analysis. We begin by $X_k \notin \partial_1$,
i.e., $X_k(1) > 0.$ There are two subcases two consider:
$k = \sigma_{j-1,j}$ for some $j \ge 1$ and 
$k \neq \sigma_{j-1,j}$ for all $j \ge 1$.
For 
$k \neq \sigma_{j-1,j}$, 
$S_k' = \Gamma_j h_{2,j,r}(T_n(X_k))$ 
$S_{k+1}' = \Gamma_j h_{2,j,r}(T_n(X_{k+1}))$ 
for some $j \in \{0,1,2,3,...,d\}$;
the functions $h_{2,j,r}$ are $Y$-superharmonic by Proposition
\ref{p:Ysuperharmonicfuns} and therefore 
$h_{2,j,r}(T_n(\cdot))$ are $X$-superharmonic on $\partial_1^c.$
It follows from these that
\begin{equation}\label{e:interiorinequalityskprime}
{\mathbb E}_x[S_{k+1}' | {\mathscr F}_k] \le S_k'
\end{equation}
over the event 
$\{ X_k \notin \partial_1\} \cap \{k \neq \sigma_{j-1,j}, j=1,2,3,...,d\}\}.$
If $k = \sigma_{j-1,j}$ for some $j \in \{2,3,4,...,d\}$ we have
$S_k' = \Gamma_{j-1}h_{2,j-1,r}(T_n(X_k))$
and $S_{k+1}' = \Gamma_j h_{2,j,r}(T_n(X_{k+1})$; $h_{2,j,r}$  is $Y$-superharmonic
and therefore $h_{2,j,r}(T_n(\cdot))$is $X$-superharmonic on $\partial_1^c$.
These imply
\begin{equation}\label{e:step1h2jr}
\Gamma_j h_{2,j,r}(T_n(X_{k}) \ge
{\mathbb E}[\Gamma_j h_{2,j,r}(T_n(X_{k+1}) |{\mathscr F}_k] =
{\mathbb E}[S_{k+1}' |{\mathscr F}_k] 
\end{equation}
over the event $\{k = \sigma_{j-1,j}\} \cap \{ X_k \notin \partial_1\}.$
That $X_{k} \in \partial_j$ for $k=\sigma_{j-1,j}$
and Lemma \ref{l:inequalityoverjumps}
imply
\[
S_{k}'=\Gamma_{j-1} h_{2,j-1,r}(T_n(X_{k}))\ 
\ge \Gamma_j h_{2,j,r}(T_n(X_k)).
\]
This and \eqref{e:step1h2jr} imply
\[
S_k' \ge {\mathbb E}[S_{k+1}' |{\mathscr F}_k] 
\]
over the event $\{k = \sigma_{j-1,j}\} \cap \{ X_k \notin \partial_1.\}.$
This and \eqref{e:interiorinequalityskprime} imply
$S_k' \ge {\mathbb E}[S_{k+1}' |{\mathscr F}_k]$;
subtracting 
$k\left( \frac{\lambda}{r} + \sum_{j=1}^d \mu_j \right) r^n$
from the left and 
$(k+1)\left( \frac{\lambda}{r} + \sum_{j=1}^d \mu_j \right) r^n$
from the right gives
\[
S_k \ge {\mathbb E}[S_{k+1} |{\mathscr F}_k] 
\]
over the event $\{X_k \notin \partial_1\}$. 

For $x \in {\mathbb Z}_+^d$, define
\[
L^*(x) \doteq \{ l \in \{2,3,4,...d\}, x(l) \neq 0 \};
\]
and 
\[
d^*(x) \doteq |L^*(x)|;
\]
$l_1(x) < l_2(x) < \cdots < l_{d^*}(x)$ are the members of $L^*$; two
conventions 1)
$l_{d^*+1}(x) = d+1$ and 2) $l_1 = d+1$
 if $d^* = 0$, i.e., if $L^*(x) = \emptyset.$
For $X_k \in \partial_1$, 
there are two cases to consider: 1) $k = \sigma_{j-1,j}$
for some $j$ and 2)  $k \neq \sigma_{j-1,j}$ for all $j$.
For the latter case 
\begin{equation}\label{e:skpp1}
S_k' = \Gamma_j h_{2,j,r}(T_n(X_k)),
S_{k+1}' = \Gamma_j h_{2,j,r}(T_n(X_{k+1})),
\end{equation}
for some $j$.
For ease of notation, let us abbreviate $l_m(X_{k})$ to $l_m$, 
and $d^*(X_{k})$ to $d^*$.
Decompose
$h_{2,j,r}(T_n(x))$ as
\begin{equation}\label{e:decomposeh2jr}
h_{2,j,r}(T_n(x)) = 
\sum_{l=1}^{(l_1-1) \wedge j}\gamma_l h_{l,r}(T_n(x)) 
+ \sum_{m=1}^{d^*} \sum_{l=l_m}^{(l_{m+1} -1)\wedge j} \gamma_l h_{l,r}(T_n(x)),
\end{equation}
where we use the conventions set above.
Let us begin by considering any of the inner sums in the
second sum in the last display. By the definition
of $l_m$ and $l_{m+1}$, $X_k(l_m) > 0$ and
$X_k(l) = 0$ for $l_m < l < l_{m+1}$. This implies
$h_{l,r}(T_n(X_k))  = h_{l_m,r}(T_n(X_k))$ 
for $l_m < l < l_{m+1}$. These, $l_m > 1$, Proposition \ref{p:hkr},
the definition \eqref{d:gamma} of $\gamma_l$ and the dynamics of $X$ imply
\begin{align*}
&{\mathbb E}
\left
[\sum_{l=l_m}^{(l_{m+1} -1)\wedge j} \gamma_l h_{l,r}(T_n(X_{k+1}))
| {\mathscr F}_k \right] -
\sum_{l=l_m}^{(l_{m+1} -1)\wedge j} \gamma_l h_{l,r}(T_n(X_k))\\
&~~~~=
h_{l,r}(T_n(X_k))
\left(
\lambda \left( \frac{1}{r} - 1 \right)
\sum_{l=l_m+1}^{(l_{m+1} -1)\wedge j} \gamma_l 
+ \gamma_{l_m} \left( \lambda\left(\frac{1}{r} -1\right) 
+ \mu_{l_m}(r-1) \right)\right)
\le 0.
\end{align*}
Summing the last inequality over $m$ gives
\begin{equation}\label{e:therestSkp}
{\mathbb E}
\left
[\sum_{m=1}^{d^*}\sum_{l=l_m}^{(l_{m+1} -1)\wedge j} \gamma_l h_{l,r}(T_n(X_{k+1}))
| {\mathscr F}_k \right] -
\sum_{m=1}^{d^*}\sum_{l=l_m}^{(l_{m+1} -1)\wedge j} \gamma_l h_{l,r}(T_n(X_k))\\
\le 0.
\end{equation}
Similarly, the definition of $l_1$ implies, $X_k(l) = 0$ for $l< l_1$,
this and $h_{l,r}(T_n(x)) = r^{n-\sum_{m=1}^l x(m)}$
imply $h_{l,r}(T_n(X_k)) = r^n$ for $l < l_1$
over the event $\{X_k \in \partial_1\}$.
These 
and the dynamics of $X$ imply
\begin{align*}
&{\mathbb E}
\left
[\sum_{l=1}^{(l_1 -1)\wedge j} \gamma_l h_{l,r}(T_n(X_{k+1}))
| {\mathscr F}_k \right] -
\sum_{l=1}^{(l_{1} -1)\wedge j} \gamma_l h_{l,r}(T_n(X_k))\\
&~~~~=
h_{1,r}(T_n(X_k))
\left(
\lambda \left( \frac{1}{r} - 1 \right)
\sum_{l=1}^{(l_{1} -1)\wedge j} \gamma_l 
\right) 
= 
r^n 
\left(
\lambda \left( \frac{1}{r} - 1 \right)
\sum_{l=1}^{(l_{1} -1)\wedge j} \gamma_l 
\right) 
\ge 0
\end{align*}
over the event $\{ X_k \in \partial_1\}.$ Putting together the last
display, \eqref{e:therestSkp}, \eqref{e:decomposeh2jr}
and $0 < \Gamma_j \le 1$
give
\begin{equation}\label{e:inforp1}
{\mathbb E}[ \Gamma_{j} h_{2,j,r}(T_n(X_{k+1}))| {\mathscr F}_k]
\le 
\Gamma_{j} h_{2,j,r}(T_n(X_{k})) + 
\left(
\lambda \left( \frac{1}{r} - 1 \right)
\sum_{l=1}^{d} \gamma_l 
\right)
\end{equation}
over the event 
$\cap_{j=1}^d \{X_k \in \partial_1, k \neq \sigma_{j-1,j}\}$;
this and \eqref{e:skpp1} imply
\[
{\mathbb E}[S_{k+1}' | {\mathscr F}_k] 
\le 
 S_{k}'  +
r^n
\left(
\lambda \left( \frac{1}{r} - 1 \right)
\sum_{l=1}^{d} \gamma_l 
\right).
\]
Moving the last expression to the left of the inequality sign
and
subtracting
$k 
\left(
\lambda \left( \frac{1}{r} - 1 \right)
\sum_{l=1}^{d} \gamma_l 
\right)
$ from both sides give
${\mathbb E}[S_{k+1} | {\mathscr F}_k] \le S_{k}$ over the same event.
It remains to show 
\begin{equation}\label{e:lastcase}
{\mathbb E}[S_{k+1} | {\mathscr F}_k] \le S_{k}
\end{equation}
over the event  
$\cup_{j=1}^d \{X_k \in \partial_1,  k = \sigma_{j-1,j} \}.$ 
In this case
\[
 S_{k}'=\Gamma_{j-1} h_{2,j-1,r}(T_n(X_{k})),
S_{k+1}' = \Gamma_j h_{2,j,r}(T_n(X_{k+1})),
\]
for some $j \in \{1,2,3,...,d\}$ and $X_k \in \partial_j.$
By Lemma \ref{l:inequalityoverjumps} 
\[
S_{k}'=\Gamma_{j-1} h_{2,j-1,r}(T_n(X_{k})) \ge
\Gamma_{j} h_{2,j,r}(T_n(X_{k}));
\]
this and \eqref{e:inforp1} imply \eqref{e:lastcase}.
\end{proof}

\end{document}
